\section{Further research questions and exact next targets}
\label{sec:research}

The results above suggest several concrete projects.  Each target is
phrased so that a successful answer has a checkable mathematical
endpoint rather than a larger table of experiments.

\begin{enumerate}[label=\textbf{R\arabic*.},leftmargin=*]
\item \textbf{Complete smaller order twelve.}
Apply the fixed-size reduction with $a=6$.  First compute a rigorous
eventual threshold and then certify the remaining finite interval with
one or several denominator-safe primes.  A stronger result would replace
the finite scan by a structural congruence or a sign argument for the
six-square determinant.

\item \textbf{Remove computation from the order-ten theorem.}
Find a recurrence, total-positivity statement, or arithmetic invariant
showing $\det\A^{[5]}\ne0$ for every $b\ge10$.  The modular vector in
this package is useful training data, but the desired endpoint is a
uniform proof explaining why cancellation cannot occur.

\item \textbf{Prove the unrestricted unequal-order conjecture.}
The eventual theorem controls $b$ for fixed $a$, but its threshold grows
with $a$.  Seek a bound on the lower Bernoulli-degree terms that is
uniform in the separation $b-a$, or a determinant factorization whose
nonvanishing is manifest for all $1\le a<b$.

\item \textbf{Closed inverse rows.}
The polynomials $P_{b,k}$ are explicit through matrix inversion but lack
a product formula.  Derive closed coefficient rows, perhaps through the
Pad\'e orthogonal polynomials and Christoffel--Darboux kernels.  A useful
answer should specify every coefficient and denominator directly, not
refer to an unevaluated inverse matrix.

\item \textbf{Minimal subtraction and accelerated identities.}
The present $Q_{b,k}$ leaves an $O(n^{-2})$ summand.  Subtract additional
negative asymptotic powers and compensate them by the known values
$\zeta(2j,1/2)$ to obtain arbitrarily rapid convergence.  Determine a
canonical acceleration with small rational coefficients and quantify
its arithmetic denominator growth.

\item \textbf{Boundary values of the two-tail generator.}
Continue \eqref{eq:two-tail-generator} through $z=1$ after the complete
Euler-operator subtraction.  Express the finite boundary values in a
minimal basis of multiple zeta values and order derivatives.  The
operator formula makes every candidate finite, but functional reduction
at high shifted order remains nontrivial.

\item \textbf{Rationally shifted and character-twisted tails.}
Replace $A_p(n)$ by $\zeta(p,n+a)$ with rational $a$, or by a Dirichlet
character tail.  Determine the analogue of the three-region numerator
in \eqref{eq:two-tail-generator}, the correct centered subtraction, and
the resulting coordinate space of Dirichlet $L$-values.

\item \textbf{Higher products.}
For three or more tails, partition the summation region by the minimum
index and derive depth-$d$ multiple-polylogarithm generators.  Then ask
for finite-rank reconstruction theorems and ordinary convergent
polygamma-product identities.  The main challenge is to organize the
ordered set partitions without hiding endpoint terms.

\item \textbf{Boundary Stieltjes correlations.}
Study the Abelian and finite-part limits of
\eqref{eq:stieltjes-correlation} as $z\to1$.  Identify exactly which
combinations of Stieltjes constants, multiple-zeta derivatives, and
powers of $\log(1-z)$ survive.  The pole-subtracted kernel should be
continued jointly in $(\alpha,\beta,z)$ rather than coefficient by
coefficient.

\item \textbf{Functional equations for the kernel.}
Investigate transformations $z\mapsto1-z$, $z\mapsto z/(z-1)$, and
roots of unity for $\K(\alpha,\beta;z)$.  A complete formula must retain
the order-derivative terms created by the pole subtractions.  Special
values could yield new exact sums of
$\gamma_r(n+1)\gamma_s(n+1)$ with algebraic weights.

\item \textbf{Higher anchored primitives and multiple Gamma functions.}
Integrate the all-index kernel repeatedly and determine when the result
reduces to classical polylogarithms, Barnes multiple Gamma functions,
or genuinely new depth.  The $r=s=0$ formula
\eqref{eq:digamma-primitive} is the first case of this hierarchy.

\item \textbf{Formal verification.}
Formalize the finite determinant reduction, the denominator-safe modular
certificate, and the min-index proof of
\eqref{eq:two-tail-generator}.  The finite-field portion is well suited
to reflection.  A later layer can formalize the analytic continuation
and centered regularization used in
\cref{thm:ordinary-reconstruction}.

\item \textbf{Independent literature and priority audit.}
Search multiple-zeta, Euler-sum, and Stieltjes-constant literature for
forms equivalent to \eqref{eq:K-kernel} under stuffle or iterated-integral
transformations.  Record equivalence maps explicitly.  This should be
completed before any publication-level novelty claim is made.
\end{enumerate}
